%% =====================================================================
%%  Prague Cafe Project -- Location topic
%%  Built from the course template. All prompts filled in from the R output tables.
%%  Tables and figures come from the R scripts in ../analysis.
%%  Lines marked "% CHECK" describe plots that were not visible when this was written.
%% =====================================================================
\documentclass[11pt,a4paper]{article}
\usepackage{cafestyle}

% Include an R-exported table if it exists, shrinking it only if it is wider than the page.
\newsavebox{\tablebox}
\newcommand{\inputtable}[1]{\IfFileExists{#1}{\sbox{\tablebox}{\input{#1}}\ifdim\wd\tablebox>\linewidth\resizebox{\linewidth}{!}{\usebox{\tablebox}}\else\usebox{\tablebox}\fi}{\par\noindent\fbox{\parbox{0.9\linewidth}{\centering Missing table file: \texttt{\detokenize{#1}}}}\par}}
\newcommand{\inputfigure}[2][0.85]{\IfFileExists{#2}{\includegraphics[width=#1\linewidth]{#2}}{\fbox{\parbox{0.8\linewidth}{\centering Missing figure file: \texttt{\detokenize{#2}}}}}}

\title{\textbf{Caf\'es Farther from the Orloj Have Fewer Ratings: Location and Popularity of Prague Caf\'es}}
\author{Tyler Bannon\\[0.3ex]\small Data Analytics, Topic: Location}
\date{\today}

\begin{document}
\maketitle

\begin{abstract}
\noindent This chapter asks where Prague caf\'es are relative to the Orloj and how that relates to how busy they are, using a student-collected sample of 189 caf\'es. Caf\'es are on average 1.94\,km from the Orloj, and they are spread about 0.42\,km farther east--west than north--south. Caf\'es farther from the Orloj have fewer Google ratings (Spearman $\rho = -0.33$), and each additional km lowers the probability of an above-median number of ratings by about 13 percentage points. Roasteries are about 0.50\,km farther out than other caf\'es, but the evidence is modest, and caf\'e types do not differ clearly. A few very popular caf\'es strongly influence the linear estimates.
\end{abstract}

\section{Introduction}
A caf\'e's popularity might depend on where it is: caf\'es close to the Old Town Square may meet more passers-by and so collect more ratings than caf\'es in outer districts. This chapter is the location chapter of the wider story about the Prague caf\'e market. It asks where caf\'es are relative to the Orloj (the Prague astronomical clock), whether the kind of caf\'e matters for location, and whether the distance from the Orloj goes together with the number of Google ratings.

\section{Data and methods}

\subsection{Data}
The data are a sample of Prague caf\'es collected by students, with the location, caf\'e type, seating and number of Google ratings entered in a shared spreadsheet. The raw file had 299 rows. The following cleaning choices were made, and each one is explained below.
\begin{itemize}
\item \textbf{99 empty rows removed}, leaving 200. They contain no information.
\item \textbf{11 caf\'es entered twice}, keeping the row entered first, leaving 189. Keeping both would count the same caf\'e twice.
\item \textbf{About 15 spellings of 9 districts harmonised.} Otherwise the same district would be split across several categories.
\item \textbf{5 caf\'es moved to the district their address shows} (sheet rows 37, 39, 43, 44, 144), because the address is more reliable than the typed district label.
\item \textbf{One blank caf\'e type set to missing} (row 136), rather than guessing a type.
\item \textbf{Paul (row 113) recoded as an international chain}, because it belongs to one.
\item \textbf{Thousands commas removed from the number of ratings}, so that values such as 23,951 are read as numbers.
\item \textbf{Caf\'e Fiesta (row 85) kept with 0 ratings}, because 0 is a possible value and dropping the caf\'e would remove a real, unpopular caf\'e from the sample.
\item \textbf{Latitude and longitude taken from the coordinates cell}, because the two disagree in row 39.
\end{itemize}
Three variables have missing values (Table~\ref{tab:desc}): the number of ratings is missing for 2 caf\'es, seats inside for 1 caf\'e, and the caf\'e type for 1 caf\'e. Each analysis uses all caf\'es with complete data for the variables it needs, so the sample size is 189 for analyses A1 to A3, 188 for A4 and A5, 187 for A6 to A10, and is reported in each table. A different choice could change the results, most of all where the groups are small: there are only 9 international chains, so recoding or dropping a single chain shifts the share of chains in Prague~1 by about 11 percentage points, and keeping a different row of a duplicated caf\'e could change a distance or a number of ratings slightly.

\subsection{Statistical methods}
All tests use a significance level of .05. For each question the test was chosen with Table~\ref{tab:decision}: a parametric test is used if a variable or group looks roughly normal or if the central limit theorem applies, and a rank-based test is used when the distribution is skewed and the smallest group is small. A one-tailed alternative is used only where the question names a direction (A2, A3, A5); all other tests are two-tailed. In the regressions A7 to A9, heteroscedasticity-consistent (HC3) standard errors are used, because the number of ratings is very skewed and a few caf\'es have large influence. Confidence intervals that are bootstrapped use 5,000 resamples, seed 123, and percentile intervals. Pairwise comparisons after the Kruskal--Wallis test in A4 use Dunn's test with Holm's adjustment. All analyses were run in R (see Appendix~\ref{app:repro}).

\section{Results}

The results begin with where the caf\'es are (A1 to A5), then turn to how many ratings they have and how that depends on distance (A6 to A10). Table~\ref{tab:summary} gives an overview of all ten analyses.

\begin{table}[ht]
\centering
\caption{Main results of all ten analyses. CI is the 95\% confidence interval, two-sided even for one-tailed tests. The statistic is $t$ for A1, A3 and A7 to A9, $W$ for A2, $H$ for A4, $S$ for A6 and $z$ for A10; Fisher's test in A5 has no test statistic. pp = percentage points.}
\label{tab:summary}
\inputtable{tables/summary.tex}
\end{table}

%% -------------------------------------------------------------------
\subsection{A1: Average distance from the Orloj}
Is the average distance of caf\'es from the Orloj different from 2\,km?

\noindent\Ho: the mean distance of caf\'es from the Orloj equals 2\,km. \Hone: the mean distance of caf\'es from the Orloj is different from 2\,km (two-tailed, fixed in advance).

\begin{table}[ht]
\centering
\caption{One-sample $t$-test of the distance from the Orloj against 2\,km. CI is the 95\% confidence interval of the mean.}
\label{tab:a1}
\inputtable{tables/A1_table.tex}
\end{table}

\begin{figure}[ht]
\centering
\inputfigure{figures/A1_figure.pdf}
\caption{Distance from the Orloj in km. Grey points are individual caf\'es, the blue point and bars are the mean and its 95\% confidence interval, and the dashed line is the 2\,km benchmark. The sample size is in Table~\ref{tab:a1}.}
\label{fig:a1}
\end{figure}

The mean distance of the 189 caf\'es from the Orloj is 1.94\,km (95\% CI 1.75 to 2.13\,km). This is not different from 2\,km, $t(188) = -0.59$, $p = .557$, and the effect size is negligible (Cohen's $d = -0.04$). The sample mean is only 0.06\,km (60 metres) below the benchmark, and a gap of this size does not matter in practice. The median is worth giving next to the mean because the distances are right-skewed (skewness 1.13, Table~\ref{tab:decision}): a few caf\'es far out, up to 8.08\,km, pull the mean up, so the typical caf\'e is closer to the Orloj, with a median of 1.73\,km and an IQR of 1.66\,km.

%% -------------------------------------------------------------------
\subsection{A2: Roasteries and other caf\'es}
Are roasteries farther from the Orloj than other caf\'es?

\noindent\Ho: the mean distance from the Orloj is the same for roasteries and other caf\'es. \Hone: the mean distance from the Orloj is greater for roasteries than for other caf\'es (one-tailed, fixed in advance). Because the descriptive analysis calls for a rank test, these are read as statements about the typical distance.

\begin{table}[ht]
\centering
\caption{Wilcoxon rank-sum test of the distance from the Orloj, roasteries minus other caf\'es. The difference is the Hodges--Lehmann estimate with its two-sided 95\% confidence interval, and the $p$-value is one-tailed and was fixed in advance.}
\label{tab:a2}
\inputtable{tables/A2_table.tex}
\end{table}

\begin{figure}[ht]
\centering
\inputfigure{figures/A2_figure.pdf}
\caption{Distance from the Orloj in km for roasteries and other caf\'es. Grey points are individual caf\'es, and the blue points and bars are the medians with bootstrapped 95\% confidence intervals (5,000 resamples, seed 123). Group sizes are in Table~\ref{tab:a2}.}
\label{fig:a2}
\end{figure}

The difference is calculated as roasteries minus other caf\'es. The 22 roasteries have a median distance of 2.17\,km and the 167 other caf\'es 1.64\,km, a Hodges--Lehmann difference of 0.50\,km (95\% CI $-0.04$ to 1.04\,km). The rank-sum test gives $W = 2274$ and a one-tailed $p = .035$. The $p$-value is one-tailed because the hypothesis names a direction, while the interval is two-sided. Together they let us claim only modest evidence that roasteries are farther out in the direction predicted: the one-tailed test passes the .05 level, but the two-sided interval just includes 0, so a difference of zero (or a slightly negative one) cannot be ruled out. Levene's test gives no evidence that the spread differs between the groups ($p = .488$). Caf\'es in the same district are not independent of each other, which makes the $p$-value somewhat too optimistic, and with only 22 roasteries the estimate is imprecise.

%% -------------------------------------------------------------------
\subsection{A3: East--west and north--south spread}
Are caf\'es spread farther east--west than north--south from the Orloj?

\noindent\Ho: the mean of the absolute east--west km minus the absolute north--south km equals 0. \Hone: this mean is greater than 0 (one-tailed, fixed in advance).

\begin{table}[ht]
\centering
\caption{Paired $t$-test of the absolute east--west distance minus the absolute north--south distance from the Orloj, in km. The interval is two-sided, and the $p$-value is one-tailed and was fixed in advance.}
\label{tab:a3}
\inputtable{tables/A3_table.tex}
\end{table}

\begin{figure}[ht]
\centering
\inputfigure{figures/A3_figure.pdf}
\caption{Within-caf\'e differences between the absolute east--west and north--south distance from the Orloj, in km. Grey points are individual caf\'es, the blue point and bars are the mean difference and its 95\% confidence interval, and the dashed line marks no difference. The sample size is in Table~\ref{tab:a3}.}
\label{fig:a3}
\end{figure}

The difference is calculated as the absolute east--west distance minus the absolute north--south distance, for each caf\'e. On average the east--west distance is larger by 0.42\,km (95\% CI 0.28 to 0.56\,km), $d_z = 0.43$, $t(188) = 5.97$, one-tailed $p < .001$. In km, the caf\'es in this sample lie about 420 metres farther from the Orloj sideways than up or down, a small to moderate effect.

%% -------------------------------------------------------------------
\subsection{A4: Distance by caf\'e type}
Does the distance from the Orloj differ by caf\'e type (independent, local chain, international chain), and if so, how?

\noindent\Ho: the mean distance from the Orloj is the same for all three caf\'e types. \Hone: the mean distance of at least one caf\'e type differs from another (no direction, fixed in advance). Because the descriptive analysis calls for a rank test, these are read as statements about the typical distance.

\begin{table}[ht]
\centering
\caption{Kruskal--Wallis test of the distance from the Orloj across the three caf\'e types. Group medians and pairwise comparisons are in Tables~\ref{tab:a4groups} and~\ref{tab:a4pairs}.}
\label{tab:a4}
\inputtable{tables/A4_table.tex}
\end{table}

\begin{figure}[ht]
\centering
\inputfigure{figures/A4_figure.pdf}
\caption{Distance from the Orloj in km by caf\'e type. Grey points are individual caf\'es, and the blue points and bars are the medians with bootstrapped 95\% confidence intervals (5,000 resamples, seed 123). Group sizes are in Table~\ref{tab:a4groups}.}
\label{fig:groups}
\end{figure}

The test gives $H(2) = 4.81$, $p = .090$ ($n = 188$), so there is no evidence at the .05 level that the distance differs by caf\'e type. The pairwise comparisons should therefore not be read as findings: they are post hoc, they follow a test that did not reject the null hypothesis, and none of the Holm-adjusted $p$-values is below .10 anyway (Appendix~\ref{app:robust}). The medians are 1.64\,km for independent caf\'es, 2.02\,km for local chains and 0.58\,km for international chains. The international chains look closer to the Orloj, but very little can be said about them: there are only 9, and their median has a wide confidence interval (0.42 to 2.36\,km) that overlaps both other groups.

%% -------------------------------------------------------------------
\subsection{A5: International chains and Prague 1}
Is an international chain more likely to be in Prague 1 than other caf\'e types, and if so, by how many percentage points?

\noindent\Ho: the share of caf\'es in Prague 1 is the same for international chains and other caf\'e types. \Hone: the share of caf\'es in Prague 1 is higher for international chains than for other caf\'e types (one-tailed, fixed in advance).

\begin{table}[ht]
\centering
\caption{Counts and row percentages of caf\'es in Prague 1 by caf\'e type.}
\label{tab:a5counts}
\inputtable{tables/A5_counts.tex}
\end{table}

\begin{table}[ht]
\centering
\caption{Fisher's exact test of international chains and Prague 1. The difference is in percentage points (international chains minus the rest). The intervals are two-sided, and the $p$-value is one-tailed and was fixed in advance.}
\label{tab:a5}
\inputtable{tables/A5_table.tex}
\end{table}

\begin{figure}[ht]
\centering
\inputfigure{figures/A5_figure.pdf}
\caption{Percent of caf\'es located in Prague 1 for international chains and other caf\'e types, with 95\% confidence intervals and the underlying counts.}
\label{fig:a5}
\end{figure}

Six of the 9 international chains (66.7\%) are in Prague 1, compared with 58 of the 179 other caf\'es (32.4\%). The difference is 34.3 percentage points (95\% CI 2.7 to 65.8 pp). The smallest expected count is 3.06, below the usual minimum of 5, so the chi-square approximation is not reliable and Fisher's exact test is used. The association is small by Cram\'er's $V$ (0.15). The odds ratio is 4.14 (95\% CI 0.85 to 26.47), and the one-tailed $p = .043$. The two intervals tell different stories: the percentage-point interval excludes 0, while the odds-ratio interval includes 1. With only nine chains, both are very wide, and a single chain in a different place would move the share by about 11 pp. More weight is given to the odds-ratio interval, which comes from the exact method and is the more cautious of the two, so the result is read as suggestive that international chains favour Prague~1, but not firmly established.

%% -------------------------------------------------------------------
\subsection{A6: Number of ratings and distance}
Is the number of Google ratings correlated with the distance from the Orloj, and if so, in which direction and how strongly?

\noindent\Ho: the population correlation between the number of ratings and distance equals 0. \Hone: it is not 0 (two-tailed, fixed in advance).

\begin{table}[ht]
\centering
\caption{Spearman correlation between the number of Google ratings and the distance from the Orloj, with a bootstrapped 95\% confidence interval (5,000 resamples, seed 123).}
\label{tab:a6}
\inputtable{tables/A6_table.tex}
\end{table}

\begin{figure}[ht]
\centering
\inputfigure{figures/A6_figure.pdf}
\caption{Number of Google ratings plus one (log scale) against the distance from the Orloj in km. Each point is one caf\'e. The title gives Spearman's $\rho$ and its bootstrapped 95\% confidence interval. The sample size is in Table~\ref{tab:a6}.}
\label{fig:a6}
\end{figure}

Spearman's rank correlation is $\rho = -0.33$ (95\% CI $-0.46$ to $-0.19$, $n = 187$), $p < .001$. Distance explains about 11\% of the variance in the ranks ($\rho^2 = 0.11$). In plain words, caf\'es that are farther from the Orloj tend to have fewer ratings, and the association is moderate in strength. This describes a pattern in the data and does not show that distance causes the difference.

%% -------------------------------------------------------------------
\subsection{A7: Ratings per km}
By how many ratings does the number of Google ratings change per km from the Orloj?

\noindent\Ho: the slope of the number of ratings on distance equals 0 ratings per km. \Hone: the slope is not 0 (two-tailed, fixed in advance).

\begin{table}[ht]
\centering
\caption{Linear regression of the number of Google ratings on the distance from the Orloj in km, with robust (HC3) standard errors and 95\% confidence intervals.}
\label{tab:a7}
\inputtable{tables/A7_table.tex}
\end{table}

\begin{figure}[ht]
\centering
\inputfigure{figures/A7_figure.pdf}
\caption{Number of Google ratings against the distance from the Orloj in km. Each point is one caf\'e, and the blue line is the regression line with its 95\% confidence band. The sample size is in Table~\ref{tab:a7}.}
\label{fig:a7}
\end{figure}

Each additional km from the Orloj goes together with 327 fewer ratings on average (slope $-327.14$, robust 95\% CI $-548.43$ to $-105.84$ ratings per km), $t = -2.90$, $p = .004$, $n = 187$. However, the model is weak: distance alone gives $R^2 = .04$. The Breusch--Pagan test gives no evidence of unequal variance in this model ($p = .189$), but the number of ratings is extremely skewed (skewness 7.83, Table~\ref{tab:decision}), so the residuals cannot be expected to be normal; this is why robust standard errors are used and why the diagnostic plots in Figure~\ref{fig:a7diag} should be read with care. % CHECK: confirm against Figure a7diag that the residuals are right-skewed and a few caf\'es lie far from the rest.
The most influential caf\'e is Caf\'e Louvre (23,951 ratings at 0.58\,km, Cook's distance 0.69), followed at a distance by Caf\'e Slavia, Caf\'e Savoy and Kav\'arna Creperi\'e (Table~\ref{tab:a7infl}). Without the four caf\'es above the $4/n$ cut-off the slope is $-196.65$ ratings per km, about 40\% smaller but still negative (Table~\ref{tab:a7sens}). A straight line in raw ratings is therefore not a good description of these data: it explains very little of the variation, and its size depends on a few very popular caf\'es. The log scale used in Figure~\ref{fig:a6} suits the data better.

%% -------------------------------------------------------------------
\subsection{A8: Accounting for seats inside}
Does the link between distance from the Orloj and the number of ratings hold once you account for the number of seats inside, and if so, how much of it remains?

\noindent\Ho: holding seats inside constant, the slope of the number of ratings on distance equals 0. \Hone: holding seats inside constant, the slope is not 0 (two-tailed, fixed in advance).

\begin{table}[ht]
\centering
\caption{Regressions of the number of Google ratings with and without seats inside, on the same caf\'es, with robust (HC3) standard errors and 95\% confidence intervals.}
\label{tab:reg}
\inputtable{tables/A8_table.tex}
\end{table}

\begin{table}[ht]
\centering
\caption{Slope of distance in ratings per km before and after adding seats inside, the percent of the slope that remains, the nested $F$-test for adding seats, and the variance inflation factor of distance.}
\label{tab:a8comp}
\inputtable{tables/A8_comparison.tex}
\end{table}

\begin{figure}[ht]
\centering
\inputfigure{figures/A8_figure.pdf}
\caption{Slope of distance in ratings per km, without and with seats inside in the model, with robust 95\% confidence intervals. The dashed line marks no association. The caf\'es behind each slope are shown in Figure~\ref{fig:a8partial}.}
\label{fig:a8}
\end{figure}

The slope of distance falls from $-327.14$ to $-229.29$ ratings per km once seats inside are added, so 70.1\% of it remains. It is still different from zero (robust 95\% CI $-408.67$ to $-49.91$, $t = -2.51$, $p = .013$). The model with seats explains more ($R^2 = .33$ against .04), and the variance inflation factor of distance is 1.01, so distance and seats are almost unrelated and collinearity is not a concern. The two tests for seats disagree: the nested $F$-test is highly significant, $F(1, 184) = 78.11$, $p < .001$, but the robust $t$-test for the seats coefficient (38.3 ratings per seat) is not, with a 95\% CI of $-16.61$ to 93.29 and $p = .173$. The reason is that the classical $F$-test assumes constant variance, and the Breusch--Pagan test rejects that for this model ($p < .001$); the robust standard error (28.04) is allowed to be larger where the variance is large. In addition, two caf\'es have very large Cook's distances (Caf\'e Louvre 8.53 and VNITROBLOCK 5.73, Table~\ref{tab:a8infl}). VNITROBLOCK is not among the influential caf\'es in A7, so its influence comes through seats inside. Without the five influential caf\'es the distance slope is $-159.40$ (Table~\ref{tab:a8sens}), smaller but still negative. The robust test is the more trustworthy of the two, so the evidence that seats matter beyond distance is weak, while distance keeps a negative association with ratings in every version of the model. % CHECK: confirm against Figure a8diag that the spread of the residuals grows with the fitted values.

%% -------------------------------------------------------------------
\subsection{A9: East and west of the Orloj}
Does the relationship between distance from the Orloj and the number of ratings differ between caf\'es east and west of the Orloj, and if so, how?

\noindent\Ho: the slope of the number of ratings on distance is the same east and west of the Orloj (the interaction is 0). \Hone: the slope differs between east and west (two-tailed, fixed in advance).

\begin{table}[ht]
\centering
\caption{Regression of the number of Google ratings on distance, side of the Orloj (1 = east) and their interaction, with robust (HC3) standard errors and 95\% confidence intervals.}
\label{tab:a9}
\inputtable{tables/A9_table.tex}
\end{table}

\begin{table}[ht]
\centering
\caption{Slope of distance in ratings per km on each side of the Orloj, calculated from the coefficients in Table~\ref{tab:a9}.}
\label{tab:a9slopes}
\inputtable{tables/A9_slopes.tex}
\end{table}

\begin{figure}[ht]
\centering
\inputfigure{figures/A9_figure.pdf}
\caption{Number of Google ratings against the distance from the Orloj in km, by side of the Orloj. Each point is one caf\'e, and the lines are regression lines with 95\% confidence bands. The sample size is in Table~\ref{tab:a9}.}
\label{fig:a9}
\end{figure}

The interaction coefficient is 477.34 ratings per km (robust 95\% CI $-231.32$ to 1186.00), $t = 1.32$, $p = .188$. The interaction term tells how much the slope of distance differs between the two sides: here the slope east of the Orloj ($-206.58$ ratings per km) is 477 ratings per km less steep than the slope west of the Orloj ($-683.92$ ratings per km). However, the confidence interval includes 0, so there is no evidence that the relationship differs between east and west; this does not show that there is no difference. The estimate is also fragile: without the three caf\'es above the $4/n$ cut-off the interaction falls to 69.49 (Table~\ref{tab:a9sens}). The west group contains the caf\'e with the most ratings in the sample (23,951), which explains why the estimate for the west is so uncertain (Table~\ref{tab:a9infl}).

%% -------------------------------------------------------------------
\subsection{A10: Probability of an above-median number of ratings}
Does the probability that a caf\'e has an above-median number of ratings change with each additional km away from the Orloj, and if so, by how many percentage points?

\noindent\Ho: the probability of an above-median number of ratings does not change with distance (the coefficient is 0). \Hone: the probability changes with distance (the coefficient is not 0; two-tailed, fixed in advance).

\begin{table}[ht]
\centering
\caption{Logistic regression of having an above-median number of Google ratings on the distance from the Orloj. The marginal effect is the average change in probability per km, in percentage points, with its 95\% confidence interval.}
\label{tab:a10}
\inputtable{tables/A10_table.tex}
\end{table}

\begin{figure}[ht]
\centering
\inputfigure{figures/A10_figure.pdf}
\caption{Probability of an above-median number of Google ratings against the distance from the Orloj in km. Grey points are individual caf\'es at 1 (above the median) or 0, moved slightly apart vertically so that they can be seen. The blue curve is the fitted logistic curve with its 95\% confidence band. The sample size is in Table~\ref{tab:a10}.}
\label{fig:a10}
\end{figure}

The outcome is 1 if a caf\'e has more ratings than the median (446, computed in code from the data) and 0 otherwise, so a caf\'e exactly on the median counts as not above it. Of the 187 caf\'es with a known number of ratings, 93 are above the median and 94 are not. Each additional km from the Orloj lowers the probability of an above-median number of ratings by 13.15 percentage points on average (95\% CI $-18.02$ to $-8.29$ pp), $z = -5.30$, $p < .001$, $n = 187$. As a check, the odds ratio is 0.55 per km, meaning the odds fall by about 45\% for each km. Because this outcome only records whether a caf\'e is above or below the median, it is not affected by how extreme the very popular caf\'es are, which makes it a more stable summary than the linear slope of A7.

%% -------------------------------------------------------------------
\section{Discussion}
Taken together, the results describe a sample in which caf\'es are on average about 1.9\,km from the Orloj and spread more east--west than north--south, and in which caf\'es farther from the Orloj tend to have fewer Google ratings. That pattern shows up in the rank correlation (A6), the regressions (A7, A8) and the logistic model (A10), where each km lowers the chance of an above-median number of ratings by about 13 percentage points. It holds, in a smaller form, when seats inside are accounted for and when the most influential caf\'es are removed. There is less to say about caf\'e types and roasteries: the evidence that roasteries are farther out is modest (A2), the types do not differ clearly in distance (A4), and the finding that international chains favour Prague~1 (A5) rests on only nine chains.

Several limits matter for how the results can be used. The data are a convenience sample collected by students, so they are not a random sample of Prague caf\'es and the results cannot be extended to all caf\'es in the city. The data are observational, so no causal claim is made: a lower number of ratings far from the Orloj could reflect fewer visitors, newer caf\'es or other differences that were not measured. Caf\'es in the same district are not independent (a third of the sample is in Prague~1, Table~\ref{tab:cat}), which means that every test assuming independence probably gives standard errors that are too small and $p$-values that are too optimistic, especially the borderline ones (A2 and A5). Finally, a few very popular caf\'es, above all Caf\'e Louvre, have a large influence on the linear estimates in A7 to A9. A reader should conclude that, in this sample, caf\'es nearer the Orloj tend to have more ratings, and should not conclude how large the effect is for Prague as a whole, why it arises, or that it differs between east and west.

\section{Conclusion}
In this sample of 189 Prague caf\'es, the typical caf\'e is about 1.7 to 1.9\,km from the Orloj, and the farther a caf\'e is from the Orloj, the fewer Google ratings it tends to have (Spearman $\rho = -0.33$; $-13$ percentage points in the probability of an above-median number of ratings per km). The size of this link depends on a few very popular caf\'es, and the evidence for differences by caf\'e type or between east and west is weak.

%% -------------------------------------------------------------------
\appendix
\clearpage

\section{Descriptive statistics}\label{app:desc}

\begin{table}[ht]
\centering
\caption{Descriptive statistics of the numeric variables used, overall and for each group compared. CI is the 95\% confidence interval of the mean. Distances are in km.}
\label{tab:desc}
\inputtable{tables/D_summary.tex}
\end{table}

\begin{table}[ht]
\centering
\caption{Counts and percentages of the categorical variables used. NA marks missing values.}
\label{tab:cat}
\inputtable{tables/D_categories.tex}
\end{table}

\clearpage
\section{Distributions and normality}\label{app:dist}
Almost every variable is right-skewed (Table~\ref{tab:decision}): distance has a skewness of 1.13 overall, the number of ratings 7.83 and seats inside 4.39; the only variable that looks roughly normal is the east--west minus north--south difference (skewness 0.41). Figures~\ref{fig:hist} to~\ref{fig:boot} show the distributions. % CHECK: confirm against the histograms and Q--Q plots that the shapes match the skewness values.
In most groups the central limit theorem applies, because there are at least 30 caf\'es and the skewness of the bootstrapped means is below 0.5, so a parametric test is justified. It does not apply to the following: roasteries (22 caf\'es), international chains (9), the number of ratings for all caf\'es (skewness of the bootstrapped means 0.62) and the number of ratings west of the Orloj (0.59). The skewness of roasteries (0.53) is very close to the limit of 0.5 but, with fewer than 30 caf\'es, the central limit theorem cannot be used for this group either, so a rank test is used. For each question, this gives: A1 (one-sample $t$-test), A2 (Wilcoxon rank-sum test), A3 (paired $t$-test), A4 (Kruskal--Wallis test), A5 (Fisher's exact test, from the expected counts), A6 (Spearman correlation) and A7 to A10 (regression models with robust standard errors, and a logistic regression).

\begin{table}[ht]
\centering
\caption{Decision table. A variable or group counts as roughly normal if its absolute skewness is below 0.5. The central limit theorem counts as applying if there are at least 30 caf\'es and the absolute skewness of 5,000 bootstrapped means (seed 123) is below 0.5. A parametric test is justified if either holds.}
\label{tab:decision}
\inputtable{tables/D_table.tex}
\end{table}

\begin{figure}[p]
\centering
\inputfigure[1]{figures/D_histograms.pdf}
\caption{Histograms of each numeric variable, overall and for each group compared.}
\label{fig:hist}
\end{figure}

\begin{figure}[p]
\centering
\inputfigure[1]{figures/D_boxplots.pdf}
\caption{Boxplots of each numeric variable, with the groups compared side by side.}
\label{fig:box}
\end{figure}

\begin{figure}[p]
\centering
\inputfigure[1]{figures/D_qqplots.pdf}
\caption{Normal Q--Q plots of each numeric variable, overall and for each group compared. Points on the red line match a normal distribution.}
\label{fig:qq}
\end{figure}

\begin{figure}[p]
\centering
\inputfigure[1]{figures/D_bootstrap.pdf}
\caption{Distribution of 5,000 bootstrapped means (seed 123) of each numeric variable, overall and for each group compared.}
\label{fig:boot}
\end{figure}

\clearpage
\section{Regression diagnostics}\label{app:diag}
The diagnostic plots for the three regressions are in Figures~\ref{fig:a7diag}, \ref{fig:a8diag} and~\ref{fig:a9diag}. In all three, Caf\'e Louvre is the most influential caf\'e (Cook's distance 0.69 in A7, 8.53 in A8 and 0.78 in A9), and in A8 VNITROBLOCK is also influential (5.73). Each conclusion survives without the caf\'es above the $4/n$ cut-off, in a weaker form. In A7 the slope falls from $-327.14$ to $-196.65$ ratings per km and in A8 from $-229.29$ to $-159.40$, so the slope stays negative. In A9 the interaction falls from 477.34 to 69.49, and it was not significant to begin with, so the conclusion of no clear difference between east and west holds. % CHECK: confirm against the diagnostic plots that the residuals are right-skewed and that the influential caf\'es stand out in the residuals-versus-leverage plots.

\begin{figure}[p]
\centering
\inputfigure[0.9]{figures/A7_diagnostics.pdf}
\caption{Diagnostic plots of the regression in A7: residuals against fitted values, normal Q--Q plot, scale--location and residuals against leverage.}
\label{fig:a7diag}
\end{figure}

\begin{table}[ht]
\centering
\caption{The five caf\'es with the largest Cook's distance in A7.}
\label{tab:a7infl}
\inputtable{tables/A7_influential.tex}
\end{table}

\begin{table}[ht]
\centering
\caption{Slope of distance in A7 with all caf\'es and without those whose Cook's distance is above $4/n$.}
\label{tab:a7sens}
\inputtable{tables/A7_sensitivity.tex}
\end{table}

\clearpage
\begin{figure}[p]
\centering
\inputfigure[0.9]{figures/A8_diagnostics.pdf}
\caption{Diagnostic plots of the regression in A8: residuals against fitted values, normal Q--Q plot, scale--location and residuals against leverage.}
\label{fig:a8diag}
\end{figure}

\begin{figure}[p]
\centering
\inputfigure[0.9]{figures/A8_partial.pdf}
\caption{Added-variable plots for A8: the number of ratings against each predictor, holding the other predictor constant. Each point is one caf\'e.}
\label{fig:a8partial}
\end{figure}

\begin{table}[ht]
\centering
\caption{The five caf\'es with the largest Cook's distance in A8.}
\label{tab:a8infl}
\inputtable{tables/A8_influential.tex}
\end{table}

\begin{table}[ht]
\centering
\caption{Slope of distance in A8 with all caf\'es and without those whose Cook's distance is above $4/n$.}
\label{tab:a8sens}
\inputtable{tables/A8_sensitivity.tex}
\end{table}

\clearpage
\begin{figure}[p]
\centering
\inputfigure[0.9]{figures/A9_diagnostics.pdf}
\caption{Diagnostic plots of the regression in A9: residuals against fitted values, normal Q--Q plot, scale--location and residuals against leverage.}
\label{fig:a9diag}
\end{figure}

\begin{table}[ht]
\centering
\caption{The five caf\'es with the largest Cook's distance in A9.}
\label{tab:a9infl}
\inputtable{tables/A9_influential.tex}
\end{table}

\begin{table}[ht]
\centering
\caption{Interaction in A9 with all caf\'es and without those whose Cook's distance is above $4/n$.}
\label{tab:a9sens}
\inputtable{tables/A9_sensitivity.tex}
\end{table}

\clearpage
\section{Post hoc tables}\label{app:robust}

\begin{table}[ht]
\centering
\caption{Median distance from the Orloj by caf\'e type, with bootstrapped 95\% confidence intervals (5,000 resamples, seed 123).}
\label{tab:a4groups}
\inputtable{tables/A4_groups.tex}
\end{table}

\begin{table}[ht]
\centering
\caption{Pairwise comparisons of caf\'e types. The difference is in medians with a bootstrapped 95\% confidence interval, and $z$ and the adjusted $p$-value come from Dunn's test with Holm's adjustment.}
\label{tab:a4pairs}
\inputtable{tables/A4_pairwise.tex}
\end{table}

\section{Reproducibility}\label{app:repro}
Run in this order from the repository root: \texttt{analysis/01\_cleaning.R}, \texttt{analysis/02\_descriptive.R}, \texttt{analysis/03\_analysis.R}. Software: R 4.6.1 with the packages dplyr, ggplot2, xtable, rstatix, car, lmtest, sandwich and marginaleffects. All bootstrap procedures use 5,000 resamples and seed 123. The random sideways spread of the points in the figures uses seed 1100.

\noindent Repository: \texttt{https://github.com/jaxsonsachs-creator/Semester\_Proj}

\end{document}
